\subsection{THE EULER INTEGRALS}

For brevity we shall say that a function \(f\) defined on an interval \(I \subset \mathbf{R}\) , and \(> 0\) on this interval, is logarithmically convex on \(I\) if \(\log f\) is convex on \(I\) . The definition of \(\Gamma(x)\) shows that this function is logarithmically convex on \(]0, +\infty[\) .

It is clear that the product of two logarithmically convex functions on I is also logarithmically convex on I. Further:

Lemma 1. Let \(f\) and \(g\) be two functions \(> 0\) and twice differentiable on an open interval I. If \(f\) and \(g\) are logarithmically convex functions on I, then \(f + g\) is logarithmically convex on I.

The relation \(\mathrm{D}^{2}\left(\log f(x)\right) \geqslant 0\) can be written \(f(x)f''(x) - \left(f'(x)\right)^{2} \geqslant 0\) . We are reduced to showing that the relations a \textgreater{} 0, \(a' > 0\) , \(ac - b^{2} \geqslant 0\) , \(a'c' - b'^{2} \geqslant 0\) imply \((a + a')(c + c') - (b + b')^{2} \geqslant 0\) ; now, when a \textgreater{} 0, the relation \(ac - b^{2} \geqslant 0\) is equivalent to the fact that the quadratic form \(ax^{2} + 2bxy + cy^{2}\) is \(\geqslant 0\) on \(R^{2}\) , and it is clear that if

\[
a x ^ {2} + 2 b x y + c y ^ {2} \geqslant 0 \quad \text { and } \quad a ^ {\prime} x ^ {2} + 2 b ^ {\prime} x y + c ^ {\prime} y ^ {2} \geqslant 0
\]

on \(\mathbf{R}^{2}\) then also \((a + a')x^{2} + 2(b + b')xy + (c + c')y^{2} \geqslant 0\) on \(\mathbf{R}^{2}\) .

Lemma 2. Let \(f\) be a finite real function, \(>0\) , defined and continuous on the product \(I \times J\) of two open intervals of \(\mathbf{R}\) , and such that, for every \(t \in J\) the function \(x \mapsto f(x,t)\) is logarithmically convex and twice differentiable on \(I\) . Under these hypotheses, if the integral \(g(x) = \int_{J} f(x,t) dt\) converges for every \(x \in I\) , then \(g\) is logarithmically convex on \(I\) .

First we show that for every compact interval \(K \subset J\) the function \(g_{K}(x) = \int_{K} f(x, t) dt\) is logarithmically convex. Indeed, if \(K = [a, b]\) , the sequence of functions

\[
g _ {n} (x) = \frac {b - a}{n} \sum_ {k = 0} ^ {n - 1} f \left(x, a + k \frac {b - a}{n}\right)
\]

converges simply to \(g_{K}(x)\) on I (II, p.~57, prop. 5) hence \(\log g_{n}\) converges simply to \(\log g_{K}\) ; by lemma 1 of VII, p.~310, \(\log g_{n}\) is convex on I, so (I, p.~27, prop. 4) it is the same for \(\log g_{K}\) .

On the other hand, g is the pointwise limit of the \(g_{K}\) along the directed set of compact subintervals of I (II, p.~64), so \(\log g\) is the pointwise limit of the \(\log g_{K}\) ; these last functions being convex on I, so is \(\log g\) (I, p.~27, prop. 4).

One can show easily that lemmas 1 and 2 remain valid even when one does not assume that the functions are twice differentiable (VII, p.~327, exerc. 5).

Lemma 3. Let \(\varphi\) be a continuous function and \(>0\) on an open interval \(J\) contained in \([0, +\infty[\) . If \(I\) is an open interval such that the integral \(g(x) = \int_{J} t^{x-1} \varphi(t) dt\) converges for all \(x \in I\) , then \(g\) is logarithmically convex on \(I\) .

Indeed, \(\log t^{x - 1} = (x - 1)\log t\) is a function of \(x\) which is convex and twice differentiable for all \(t > 0\) , so lemma 2 applies.

PROPOSITION 3. For all \(x > 0\)

\[
\Gamma (x) = \int_ {0} ^ {\infty} e ^ {- t} t ^ {x - 1} d t\tag{14}
\]

(second Euler integral).

Now the function \(g(x) = \int_0^\infty e^{-t} t^{x-1} dt\) is defined for all \(x > 0\) (V, p.~229); lemma 3 of VII, p.~311, then shows that it is logarithmically convex on \(]0, +\infty[\) . Moreover, on integrating by parts, one has

\[
g (x + 1) = \int_ {0} ^ {\infty} e ^ {- t} t ^ {x} d t = - e ^ {- t} t ^ {x} \Big | _ {0} ^ {\infty} + x \int_ {0} ^ {\infty} e ^ {- t} t ^ {x - 1} d t = x g (x).
\]

In other words, g is a solution of equation (1) of VII, p.~305; finally,

\[
g (1) = \int_ {0} ^ {\infty} e ^ {- t} d t = 1;
\]

the proposition therefore follows from prop. 1 of VII, p.~306.

By the change of variable \(e^{-t}=u\) one deduces from (14) (VII, p.~311) the formula

\[
\Gamma (x) = \int_ {0} ^ {1} \left(\log \frac {1}{t}\right) ^ {x - 1} d t.\tag{15}
\]

Similarly, from the change of variable \(u = t^{x}\) we obtain

\[
x \Gamma (x) = \int_ {0} ^ {\infty} e ^ {- t ^ {1 / x}} d t
\]

or again, taking account of (7) (VII, p.~3),

\[
\Gamma \left(1 + \frac {1}{x}\right) = \int_ {0} ^ {\infty} e ^ {- t ^ {x}} d t\tag{16}
\]

and in particular, for \(x = 2\)

\[
\Gamma \left(\frac {3}{2}\right) = \frac {1}{2} \Gamma \left(\frac {1}{2}\right) = \int_ {0} ^ {\infty} e ^ {- t ^ {2}} d t.\tag{17}
\]

PROPOSITION 4. For x \textgreater{} 0 and y \textgreater{} 0 the integral

\[
\mathbf {B} (x, y) = \int_ {0} ^ {1} t ^ {x - 1} (1 - y) ^ {y - 1} d t
\]

(first Euler integral) has the value

\[
\mathbf {B} (x, y) = \frac {\Gamma (x) \Gamma (y)}{\Gamma (x + y)}.\tag{18}
\]

Indeed this integral converges for \(x > 0\) and \(y > 0\) (V, p.~229). By lemma 3 of VII, p.~311, the function \(x \mapsto \mathbf{B}(x, y)\) is logarithmically convex for \(x > 0\) . Moreover,

\[
\mathbf {B} (x + 1, y) = \int_ {0} ^ {1} (1 - t) ^ {x + y - 1} \left(\frac {t}{1 - t}\right) ^ {x} d t
\]

whence, on integrating by parts,

\[
\begin{array}{l} \mathbf {B} (x + 1, y) = - \frac {(1 - t) ^ {x + y}}{x + y} \left(\frac {t}{1 - t}\right) ^ {x} \Bigg | _ {0} ^ {1} \\ \qquad + \frac {x}{x + y} \int_ {0} ^ {1} (1 - t) ^ {x + y} \left(\frac {t}{1 - t}\right) ^ {x - 1} \frac {d t}{(1 - t) ^ {2}} = \frac {x}{x + y} \mathbf {B} (x, y). \end{array}
\]

It follows that \(f(x) = \mathbf{B}(x, y) \Gamma(x + y)\) satisfies the identity (1) of VII, p.~305. Moreover, this function is logarithmically convex, being the product of two logarithmically convex functions. Finally, one has \(f(1) = \mathbf{B}(1, y) \Gamma(y + 1)\) , and \(\mathbf{B}(1, y) = \int_{0}^{1} (1 - t)^{y - 1} dt = 1/y\) , whence \(f(1) = \frac{1}{y} \Gamma(y + 1) = \Gamma(y)\) . The function \(f(x)/\Gamma(y)\) is thus equal to \(\Gamma(x)\) by prop. 1 of VII, p.~306, which proves (18).

By the change of variable \(t = \frac{u}{u + 1}\) the formula (18) becomes

\[
\int_ {0} ^ {\infty} \frac {t ^ {x - 1}}{(1 + t) ^ {x + y}} d t = \frac {\Gamma (x) \Gamma (y)}{\Gamma (x + y)}\tag{19}
\]

and by the change of variable \(t = \sin^{2} \varphi\)

\[
\int_ {0} ^ {\pi / 2} \sin^ {2 x - 1} \varphi \cos^ {2 y - 1} \varphi d \varphi = \frac {1}{2} \frac {\Gamma (x) \Gamma (y)}{\Gamma (x + y)}.\tag{20}
\]

If one puts \(x = y = \frac{1}{2}\) in this last formula it follows that

\[
\Gamma (\frac {1}{2}) = \sqrt {\pi}\tag{21}
\]

whence, by (17),

\[
\int_ {0} ^ {\infty} e ^ {- t ^ {2}} d t = \frac {1}{2} \sqrt {\pi}.\tag{22}
\]

From the relation (7) of VII, p.~307, one has the asymptotic expansion

\[
\begin{array}{l} \Gamma (x) = \frac {1}{x} \Gamma (x + 1) \\ = \frac {1}{x} + \Gamma^ {\prime} (1) + \frac {1}{2 !} \Gamma^ {\prime \prime} (1) x + \dots + \frac {1}{n !} \Gamma^ {(n)} (1) x ^ {n - 1} + O (x ^ {n}) \end{array}\tag{23}
\]

for \(\Gamma(x)\) , on a neighbourhood of 0.

Similarly, for all y fixed and \textgreater{} 0 one can write

\[
\begin{array}{l} \frac {1}{\Gamma (x + y)} = \frac {1}{\Gamma (y)} + D \left(\frac {1}{\Gamma (y)}\right) x \\ \qquad + \frac {1}{2 !} D ^ {2} \left(\frac {1}{\Gamma (y)}\right) x ^ {2} + \dots + \frac {1}{n !} D ^ {n} \left(\frac {1}{\Gamma (y)}\right) x ^ {n} + O _ {1} (x ^ {n + 1}) \end{array}
\]

and the formula (18) then gives, for y fixed, the asymptotic expansion

\[
\begin{array}{l} \mathbf {B} (x, y) = \frac {1}{x} + \left(\Gamma^ {\prime} (1) - \frac {\Gamma^ {\prime} (y)}{\Gamma (y)}\right) \\ \qquad + \left(\frac {\Gamma^ {\prime \prime} (1)}{2} - \Gamma^ {\prime} (1) \frac {\Gamma^ {\prime} (y)}{\Gamma (y)} + \frac {2 \Gamma^ {\prime 2} (y) - \Gamma (y) \Gamma^ {\prime \prime} (y)}{2 \Gamma^ {2} (y)}\right) x + O (x ^ {2}) \end{array}\tag{24}
\]

on a neighbourhood of \(x = 0\) .

Moreover, for x \textgreater{} 0 and y \textgreater{} 0 one has

\[
\begin{array}{l} \mathbf {B} (x, y) = \int_ {0} ^ {1} \left(t ^ {x - 1} + t ^ {x} \frac {(1 - t) ^ {y - 1} - 1}{t}\right) d t \\ = \frac {1}{x} + \int_ {0} ^ {1} t ^ {x} \frac {(1 - t) ^ {y - 1} - 1}{t} d t. \end{array}\tag{25}
\]

The function \(\varphi(t) = \frac{(1 - t)^{y - 1} - 1}{t}\) is continuous on the compact interval {[}0, 1{]}; since

\[
t ^ {x} = e ^ {x \log t} = 1 + x \log t + \frac {x ^ {2}}{2 !} (\log t) ^ {2} + \dots + \frac {x ^ {n}}{n !} (\log t) ^ {n} + r _ {n} (x, t)
\]

with \(|r_n(x,t)| \leqslant \frac{x^{n+1}}{(n+1)!} |\log t|^{n+1}\) (since \(\log t \leqslant 0\) and \(x > 0\) ), formula (25) gives the asymptotic expansion

\[
\begin{array}{l} \mathbf {B} (x, y) = \frac {1}{x} + \int_ {0} ^ {1} \varphi (t) d t + x \int_ {0} ^ {1} \varphi (t) \log t d t + \dots \\ \qquad + \frac {x ^ {n}}{n !} \int_ {0} ^ {1} \varphi (t) (\log t) ^ {n} d t + O _ {2} (x ^ {n + 1}) \end{array}
\]

for \(\mathbf{B}(x,y)\) on a neighbourhood of 0.

For \(n = 1\) the identification of this expansion with (24) gives in particular

\[
\Gamma^ {\prime} (1) - \frac {\Gamma^ {\prime} (y)}{\Gamma (y)} = \int_ {0} ^ {1} \frac {(1 - t) ^ {y - 1} - 1}{t} d t.
\]

Furthermore, the formula (10) gives \(\Gamma'(1) = \Gamma'(1)/\Gamma(1) = -\gamma\) , so (Gauss' integral)

\[
\frac {\Gamma^ {\prime} (x)}{\Gamma (x)} + \gamma = \int_ {0} ^ {1} \frac {1 - (1 - t) ^ {x - 1}}{t} d t.\tag{26}
\]

